% ---
% id: protocol-compiler/def:weighted
% kind: def
% title: "Choice-weighted families"
% summary: >-
%   A public choice enters as a weight on a complete conditional measurement:
%   one declared weight per choice value, or, at a measured factor after its
%   chain, a stochastic weight read through its station's setting. Weights and
%   families are binary64 numbers read exactly; products are formed exactly.
% prominence: standard
% status: implemented
% uses: [def:compile, def:exact, def:declaration]
% mentions: [lem:weighted, rem:relabel, thm:choice-encoding, ass:coin, def:spec, def:chaincond, rem:condfamily, ex:condfamily, reductions/def:joint-family, foundations/req:declared-order, req:derived-objects]
% assumes: []
% refs: []
% provenance: newly-derived
% implements: certiqs_key.protocol_compiler.weights:decide_conditional_families
% lean:
% notes: "The alphabets T and the choice alphabet are written inline as \\mathcal{T} and \\mathcal{F}; the macros \\Tset and \\Fset of spec/macros.tex are the optical instance's alphabets and are not used for the generic setting and choice."
% ---

A per-round public choice of known bias --- a basis choice, a choice between
test and generation rounds, a probe setting --- is encoded as a weight inside
the measurement operators and never as a separate fraction. There are two
constructions: one declared weight per choice value, independent of every
setting, and, at a \tref{measured factor} after the \tref{reduction chain}, a weight
that depends on the private setting of the station that measures it.

Assumption~\ref{ass:coin} and Remark~\ref{rem:relabel} are why a
setting-dependent weight is read through the station's setting and through
nothing else, and Remark~\ref{rem:relabel} and Theorem~\ref{thm:choice-encoding} why
that form is sound after the chain and is not a link of it. Which conditions
of the weighted family do and do not follow from the decision on the
conditional families, and why that decision is taken here rather than among
the target conditions, is Remark~\ref{rem:condfamily}; that the weighted
family can satisfy all five while none of the conditional families it was
built from is a POVM is Example~\ref{ex:condfamily}. The order in which the
weights are summed is declared with the public-choice record, as
\xref{foundations/req:declared-order} requires of every summed predicate, and
a presentation indexed by a product alphabet is refused under
Requirement~\ref{req:derived-objects}.

\begin{definition}\label{def:weighted}
Let $S$ be a finite-dimensional system, $\Tpub$ a finite non-empty
\term[The finite alphabet of values of a per-round public choice of known bias, each value with a declared positive weight.]{public-choice alphabet}
with declared weights $q_t>0$, $\sum_{t\in\Tpub}q_t=1$, and for each
$t\in\Tpub$ a POVM $\{N_{u|t}\}_{u\in\Uset_t}$ on $S$. The
\term[The measurement whose elements are declared measurement elements multiplied by the declared weight of a public choice value.]{choice-weighted family}
is indexed by the disjoint union
$\Vset\coloneqq\bigsqcup_{t\in\Tpub}\Uset_t\times\{t\}$ and defined by
\begin{equation}\label{eq:weighted}
  N_{u,t}\coloneqq q_t\,N_{u|t},\qquad (u,t)\in\Vset .
\end{equation}
Several independent choices compose by iterating \eqref{eq:weighted}; the
resulting weight is the product of the individual weights. \emph{What is
declared, and therefore what is decided, is the innermost family and each
weight list}: an iterated presentation declares $\{N_{u|t}\}_u$ for each value
of each choice together with one weight list per choice, the decision on the
conditional families is taken on those, and the intermediate
families are this module's own products, which inherit admissibility by
Lemma~\ref{lem:weighted}\,(i) at each step and carry no second declared
condition; the representability decision is taken on every product the map
forms, the intermediate ones included. This is the
\term[The choice-weighted family whose weights are one declared number per choice value, the same in every round.]{setting-independent form}:
the weight of a choice value is one declared number, the same in every round
and whatever else either party's apparatus did.

\emph{A weighting that depends on a station's setting.} Let
$\{F_v\}_{v\in\Vset}$ be a POVM on $S$ over a finite alphabet $\Vset$ with a
\tref{declared order} --- the \tref{target family} of the last link of $\Chrec$, as
Definition~\ref{def:compile} takes it --- and let
$t:\Vset\to\mathcal{T}$ be a declared total map onto a finite non-empty
\term[The finite alphabet of the private per-round setting of the station that measures the family, onto which the declared setting map sends each outcome.]{setting alphabet},
the
\term[The value of the measuring station's private setting that an outcome records, read by the declared setting map.]{setting component}
of the outcome: the value of the private per-round setting of the measuring
apparatus that the outcome $v$ records. Where the chain was emitted by the reductions module, $t$ is the
class map of that module's \tref{joint alphabet}
(\xref{reductions/def:joint-family}\,(a): $\Vset=\mathcal{Y}\times\mathcal{T}$
with classes $\Vset_t$); where the family is declared directly, $t$ is
declared with it, the singleton $\mathcal{T}=\{\ast\}$ being the case of no
setting. Let $\mathcal{F}$ be a finite non-empty
\term[The finite alphabet, with a declared order, of a public choice whose weight depends on the private setting of the station that measures the family.]{choice alphabet}
with a declared order and let $P(f\,|\,t)$,
$(t,f)\in\mathcal{T}\times\mathcal{F}$, be a \emph{stochastic matrix}:
$P(f\,|\,t)\ge0$ and $\sum_{f\in\mathcal{F}}P(f\,|\,t)=1$ for every
$t\in\mathcal{T}$. The
\term[The pairs of an outcome and a choice value whose weight under the declared stochastic matrix is positive.]{support of the weighting}
is
\begin{equation}\label{eq:support}
  \Vset_P\coloneqq\big\{(v,f)\in\Vset\times\mathcal{F}:\ P(f\,|\,t(v))>0\big\},
\end{equation}
ordered by the product order --- $v$ in the declared order of $\Vset$ and,
within each $v$, $f$ in the declared order of $\mathcal{F}$, the pairs off
the support omitted --- and the choice-weighted family in its
\term[The choice-weighted family of a measured system after its chain: each delivered element times a stochastic weight read through the measuring station's setting.]{setting-dependent form}
is
\begin{equation}\label{eq:weighted-t}
  N_{v,f}\coloneqq P(f\,|\,t(v))\,F_v,\qquad (v,f)\in\Vset_P .
\end{equation}
The setting-independent form is the case $P(f\,|\,t)=q_f$ for every $t$, in
which \eqref{eq:weighted-t} is \eqref{eq:weighted} with
$\Tpub=\mathcal{F}$ and $N_{v|f}=F_v$ for every $f$ --- one conditional
family, weighted by a choice that does not see the setting; conversely,
\eqref{eq:weighted} with conditional families that differ with the choice is
not of the form \eqref{eq:weighted-t}, and the two forms are two constructions
overlapping in that case. \emph{In the overlap the two constructions do not
agree on an order, and the order is a property of the presentation and not of
the family}: the alphabet the first form produces is
$\bigsqcup_f\Uset_f\times\{f\}$, listed choice-major, while the support
\eqref{eq:support} of the second is listed outcome-major, so the two are
transposes of one another and one family declared through the two records has
its condition~(S) decided in two different orders. (S) is order-dependent by
construction (Definition~\ref{def:exact}), and the two orders are not
reconciled: what fixes the order is the record that was declared, and a
\tref{declarer} that declares the same family both ways must expect the
decision its own record's order gives. \emph{$P$ is read through the setting
and through nothing else.} It is declared on $\mathcal{T}\times\mathcal{F}$
and reaches an outcome only through $t(v)$. Zero entries of $P$ are admitted, and the support
is what they are for: a pair $(v,f)$ with $P(f\,|\,t(v))=0$ is not an outcome,
is not emitted and indexes nothing, so every emitted element of
\eqref{eq:weighted-t} carries a strictly positive weight, exactly as $q_t>0$
requires of every element of \eqref{eq:weighted}. Which choice values a given
setting admits is therefore declared by the zero pattern of $P$ and not by a
zero-weight outcome. The
form is applied by the compiler \emph{after} the reduction chain, to the
family the last link delivered (Definition~\ref{def:compile}).

\emph{The reading of each declared condition.} The weights are
\tref{binary64 numbers} and the conditional families binary64 matrices. So
$\sum_tq_t=1$ is decided
\emph{exactly}, with the sum accumulated in the declared order of $\Tpub$ ---
the order in which the declared public-choice record lists $\Tpub$, fixed
there because $\idxm$ does not index $\Tpub$
(Definition~\ref{def:compile}) --- and
required to \tref[rounds nowhere]{round nowhere}; and ``$\{N_{u|t}\}_u$ is a
POVM on $S$'' is exact admissibility in the sense of
Definition~\ref{def:exact}, per block where $S$ carries sectors: (H), (P) and
(C) per block; (S) on the completeness sum $\sum_{u\in\Uset_t}N_{u|t}$,
accumulated in the order stated below and required to round nowhere; and (T)
as the agreement --- in block kind, in block identifiers, in their order and in
their dimensions --- of the maps \emph{the conditional families of one
presentation themselves carry}: there is one such map for the presentation,
and the family the weighting produces is a \tref{block list} against
\emph{it}. On a system carrying no sector decomposition, Alice's $A$, (T) is
empty. This condition is decided of each $t\in\Tpub$ separately, on the
declared numbers.

\emph{In which order (S) is read on a conditional family.} The family the
weighting produces is indexed by one alphabet, and that alphabet carries one
order: $\Xset$ where the presentation is Alice's, which arrives with its
declared order (Definition~\ref{def:spec}\,(a)); the physical receiver
alphabet where it is a measuring party's, which is the order the fibre sum of
Definition~\ref{def:chaincond}\,(C4) is accumulated in and, in \tref{Band A}, is
the emitted alphabet's own. Condition~(S) on the conditional family
$\{N_{u|t}\}_u$ is decided in the \emph{restriction of that order} to the
elements the choice value $t$ contributes --- a total order on $\Uset_t$
whatever the presentation listed, so no compatibility of the record's listing
with it is required and none is decided. Where a declared order arrives with
that alphabet the compiler takes it, by the binding rule of
Definition~\ref{def:compile}, to which this definition adds nothing; where
none arrives, the choice the compiler makes and emits is the record's own
listing, $t$ in the order it lists $\Tpub$ in and, within each $t$, $u$ in the
order it lists $\Uset_t$ in. \emph{No presentation is refused for an order}:
what is fixed here is which order a decision is read in, not what a declarer
may declare.

\emph{The remaining declared conditions, at the same reading.} The strict
inequality $q_t>0$ is a domain condition on a declared real scalar, decided by
its sign, and Definition~\ref{def:exact} carries no
strict-inequality row because none is wanted --- $q_t\ne0$ is what
Lemma~\ref{lem:weighted}\,(ii) needs, and a negative $q_t$ would reverse the
sign of an element of the emitted family, which no public choice of known bias
can mean. For the setting-dependent form the same reading applies to $P$: each
entry is a binary64 number, $P(f\,|\,t)\ge0$ is decided by sign, each row sum
$\sum_fP(f\,|\,t)=1$ is decided exactly in the declared order of
$\mathcal{F}$ and required to round nowhere, the support \eqref{eq:support}
is decided by the sign of each entry, and ``$\{F_v\}$ is a POVM on $S$'' is
the exact admissibility the chain's (C3) already decided of the family the
last link delivered. \emph{This form has no second conditional family to
decide}: the choice value is drawn after the measurement, $\{F_v\}$ is the
whole of what the decision on the conditional families has as its subject
here, and (C3) is where it was taken.

\emph{What this module constructs, and the two emission decisions.} The
products $q_tN_{u|t}$ of \eqref{eq:weighted} and $P(f\,|\,t(v))F_v$ of
\eqref{eq:weighted-t} are formed \emph{exactly} on the declared dyadic
numbers. A product of two binary64 numbers need not be one, so the compiler
decides, before emitting anything, that every entry of every $N_{u,t}$,
respectively of every $N_{v,f}$, is a binary64 number --- on every product
this map forms, whether the product is emitted or handed on --- and that the
emitted family satisfies Definition~\ref{def:exact} --- (H), (P), (C) and, on
its own declared-order sum over $\Vset$, respectively over $\Vset_P$ in its
product order, (S). These are the two \emph{emission decisions}.

\emph{The lattice of the products.} Let the arriving entries be integer
multiples of $2^{-k}$ of modulus at most $1$, and let each $q_t$ be an
integer multiple of $2^{-k_q}$; then every entry of $q_tN_{u|t}$ is an integer
multiple of $2^{-(k+k_q)}$ of modulus at most $1$, so the constructed family
lies on a lattice $k_q$ steps finer than the one it was handed and is exactly
binary64 whenever $k+k_q\le53$. Where a $q_t$ is not a power of two the finer
lattice can leave binary64 altogether: the decimal literal $0.9$ is an integer
multiple of $2^{-53}$ and of no coarser one, so $k+k_q>53$ for every $k\ge1$
and the exact product of a declared entry by it need not be a binary64
number. The same arithmetic governs $P(f\,|\,t(v))F_v$ with the exponent
$k_P$ of the entries of $P$ in place of $k_q$: a $P$ whose non-zero entries
are dyadic with small exponent --- $\tfrac12$, $\tfrac14$, $1$ --- composes
with any family the chain delivers at $k\le53-k_P$, and a decimal literal in
$P$ can leave binary64 exactly as a decimal $q_t$ can.

\emph{The general format.} In the general format of
Definition~\ref{def:declaration}\,(DX1)--(DX6) both forms above are read
\emph{on the system of the party that declares them}, and nothing else of
this definition changes. At one \tref{measured factor}, in \tref{POVM form}
and with no factor tagged $A$, the readings below are the text above,
sentence for sentence.

\emph{Which system $S$ is.} For a presentation of Alice's it is $A$, and (T)
is \emph{not empty where a factor is tagged $A$}: the conditional families
are then block lists along $\mathbb{T}$ and (T) is the four-component
agreement of the maps they carry with that factor's map, as on any other
sectored system. For a presentation of a measuring party's it is that party's
\emph{own} measured factor: one presentation per factor, decided against that
factor's own map; \emph{a presentation indexed by a product alphabet has no
declared system to be a family on} --- the product is the rule
\eqref{eq:productfamily} and is no datum (Definition~\ref{def:compile}, the
general format). So the decision ``$\{N_{u|t}\}_u$ is
a POVM on $S$'' and the products this module forms are per factor, and the
representability decision keeps exactly the subject it has: the products
\emph{these two constructions} form, and no product block of
\eqref{eq:productfamily}.

\emph{The \tref{ensemble form} of Alice's slot.} Where Alice's measurement slot
is declared in ensemble form (DX4), a choice-weighted presentation of it
declares, for each $t\in\Tpub$, a family $\{N_{x|t}\}_{x\in\Xset}$ of block
lists along $\mathbb{T}$ with every block Hermitian and positive
semidefinite, and \emph{in place of the POVM condition} the ensemble
condition of \eqref{eq:ensembleform}: for every Alice sector $t'$,
$\sum_{x}N^{(t')}_{x|t}=\tau^{(t')}_A$, the sum accumulated in the order
stated above and rounding nowhere. That condition is decided \emph{of each
$t$ separately and on the declared numbers}, with (H), (P) per block, that
identity in place of (C), and (T) as above. \emph{Why the
construction is the same construction.} With $N_{x,t}\coloneqq q_tN_{x|t}$ of
\eqref{eq:weighted}, the exact real numbers the data denote satisfy
\begin{equation}\label{eq:weighted-ensemble}
  \sum_{(x,t)\in\Vset}N^{(t')}_{x,t}
  =\sum_{t\in\Tpub}q_t\sum_{x\in\Xset}N^{(t')}_{x|t}
  =\Big(\sum_{t\in\Tpub}q_t\Big)\tau^{(t')}_A=\tau^{(t')}_A
\end{equation}
for every $t'$, by the decision just taken and by $\sum_tq_t=1$; so a
choice-weighted presentation in ensemble form whose conditional ensembles are
admitted denotes a weighted family with the ensemble identity of
\eqref{eq:ensembleform}. \emph{That is a consistency statement and not a
substitute for the decision on the emitted family}: the emitted family's own
(S), (H), (P) and its identity with the declared $\tau_A$ are decided at check
V3, on the emitted blocks and in the emitted order, and neither condition
implies the other: \eqref{eq:weighted-ensemble} is an identity of real
numbers and the emitted decision is about the additions that compute it. In
\tref{untrusted-source mode} there is no ensemble form and this paragraph has
no subject.

\emph{The setting-dependent form, per measured factor.} It is applied by the
compiler after \emph{that factor's} \tref{receiver segment}, to the family that
factor's last link delivered, with that factor's own setting map, support
\eqref{eq:support} and product order; the setting component read is a
component of \emph{that factor's} outcome, the private per-round setting of
the station that measures that factor, whichever party's label the factor
carries. This form has no second
conditional family here either, per factor: (C3) of
Definition~\ref{def:chaincond}, read for that factor, is where the delivered
family was decided.
\end{definition}
